\documentclass[11pt]{article}
\usepackage[margin=1in]{geometry}
\usepackage{amsmath, amssymb, amsthm}

\newtheorem*{claim}{Claim}
\newtheorem*{lemma}{Lemma}

\newcommand{\R}{\mathbb{R}}
\newcommand{\N}{\mathbb{N}}

\title{Q2: The Convex Hull of a Finite Set Is Compact}
\date{}

\begin{document}
\maketitle

\section*{Problem}

Fix $k, n \in \N^{+}$. Let $S = \{p_1, \dots, p_k\}$ be a finite set of points of $\R^n$.
Define the convex hull of $S$ to be the set
\[
  C = \left\{ \lambda_1 p_1 + \dots + \lambda_k p_k : \lambda_1 \ge 0, \dots, \lambda_k \ge 0,\ \lambda_1 + \dots + \lambda_k = 1 \right\} \subseteq \R^n .
\]
Prove that $C$ is compact using topological properties of continuity.

\section*{Strategy}

Write $C$ as the image of a compact set under a continuous map:
\begin{enumerate}
  \item Define the standard simplex $\Delta \subseteq \R^k$ of admissible coefficient vectors.
  \item Define a continuous map $f \colon \R^k \to \R^n$ with $f(\Delta) = C$.
  \item Show $\Delta$ is compact (closed via preimages of closed sets, and bounded).
  \item Conclude, since the continuous image of a compact set is compact (proved with sequences).
\end{enumerate}

\section*{Solution}

\paragraph{Step 1: The simplex and the map.}
Let
\[
  \Delta = \left\{ \lambda = (\lambda_1, \dots, \lambda_k) \in \R^k : \lambda_i \ge 0 \text{ for all } i,\ \sum_{i=1}^k \lambda_i = 1 \right\},
\]
and define $f \colon \R^k \to \R^n$ by
\[
  f(\lambda) = \sum_{i=1}^k \lambda_i p_i .
\]
By the definition of $C$, we have exactly $C = f(\Delta)$.

\paragraph{Step 2: $f$ is continuous.}
Let $M = \max_{1 \le i \le k} \|p_i\|$. For each $i$,
\[
  |\lambda_i - \mu_i| = \left( (\lambda_i - \mu_i)^2 \right)^{1/2}
  \le \left( \sum_{j=1}^k (\lambda_j - \mu_j)^2 \right)^{1/2} = \|\lambda - \mu\| .
\]
So for $\lambda, \mu \in \R^k$, by the triangle inequality,
\[
  \|f(\lambda) - f(\mu)\|
  = \left\| \sum_{i=1}^k (\lambda_i - \mu_i) p_i \right\|
  \le \sum_{i=1}^k |\lambda_i - \mu_i| \, \|p_i\|
  \le \sum_{i=1}^k M \|\lambda - \mu\|
  = kM\, \|\lambda - \mu\| .
\]
Now fix $\mu \in \R^k$ and $\varepsilon > 0$, and set $\delta = \dfrac{\varepsilon}{kM + 1} > 0$.
If $\|\lambda - \mu\| < \delta$, then
\[
  \|f(\lambda) - f(\mu)\| \le kM\, \|\lambda - \mu\| < kM\, \delta
  = \frac{kM}{kM + 1}\, \varepsilon \le \varepsilon .
\]
So $f$ is continuous at every $\mu \in \R^k$, i.e.\ $f$ is continuous.
(The $+1$ in the denominator only ensures $\delta$ is defined when $M = 0$.)

\paragraph{Step 3: $\Delta$ is closed.}
For each $i$, the coordinate projection $\pi_i \colon \R^k \to \R$, $\pi_i(\lambda) = \lambda_i$,
is continuous: given $\varepsilon > 0$, take $\delta = \varepsilon$; then $\|\lambda - \mu\| < \delta$
implies $|\pi_i(\lambda) - \pi_i(\mu)| = |\lambda_i - \mu_i| \le \|\lambda - \mu\| < \varepsilon$.
The sum map $s \colon \R^k \to \R$, $s(\lambda) = \sum_{i=1}^k \lambda_i = \sum_{i=1}^k \pi_i(\lambda)$,
is a finite sum of continuous functions, hence continuous. Therefore
\[
  \Delta = \left( \bigcap_{i=1}^k \pi_i^{-1}\big([0, \infty)\big) \right) \cap s^{-1}\big(\{1\}\big).
\]
Since $[0, \infty)$ and $\{1\}$ are closed in $\R$, and the preimage of a closed set under a
continuous map is closed, each set in this intersection is closed in $\R^k$. An intersection of
closed sets is closed, so $\Delta$ is closed.

\paragraph{Step 4: $\Delta$ is bounded.}
If $\lambda \in \Delta$, then each $\lambda_i \ge 0$ and $\lambda_i \le \sum_{j=1}^k \lambda_j = 1$,
so $0 \le \lambda_i \le 1$. Hence $\Delta \subseteq [0,1]^k$ and
\[
  \|\lambda\| = \left( \sum_{i=1}^k \lambda_i^2 \right)^{1/2} \le \sqrt{k} \quad \text{for all } \lambda \in \Delta .
\]
So $\Delta$ is bounded.

\paragraph{Step 5: $\Delta$ is compact.}
Recall that $K \subseteq \R^m$ is compact if every sequence in $K$ has a subsequence converging
to a point of $K$, and that this is equivalent to $K$ being closed and bounded. By Steps 3 and 4,
$\Delta$ is closed and bounded, so $\Delta$ is compact.

\paragraph{Step 6: The continuous image of a compact set is compact.}
\begin{lemma}
  If $g \colon \R^k \to \R^n$ is continuous and $K \subseteq \R^k$ is compact, then $g(K)$ is compact.
\end{lemma}
\begin{proof}
  Let $(y_m)_{m \ge 1}$ be a sequence in $g(K)$. For each $m$ choose $x_m \in K$ with
  $g(x_m) = y_m$. Since $K$ is compact, there is a subsequence $(x_{m_j})_{j \ge 1}$ and a point
  $x \in K$ with $x_{m_j} \to x$.

  We claim $y_{m_j} = g(x_{m_j}) \to g(x)$. Let $\varepsilon > 0$. By continuity of $g$ at $x$,
  there is $\delta > 0$ such that $\|z - x\| < \delta$ implies $\|g(z) - g(x)\| < \varepsilon$.
  Since $x_{m_j} \to x$, there is $J$ such that $\|x_{m_j} - x\| < \delta$ for all $j \ge J$.
  Hence $\|g(x_{m_j}) - g(x)\| < \varepsilon$ for all $j \ge J$, proving the claim.

  So $(y_m)$ has a subsequence converging to $g(x)$, and $g(x) \in g(K)$ because $x \in K$.
  Therefore $g(K)$ is compact.
\end{proof}

\paragraph{Conclusion.}
By Step 2, $f$ is continuous; by Step 5, $\Delta$ is compact. Applying the Lemma with
$g = f$ and $K = \Delta$,
\[
  C = f(\Delta)
\]
is compact. \qed

\end{document}
